\documentclass[10pt,a4paper]{amsart}
\usepackage[margin=19mm]{geometry}
\usepackage{amsmath,amssymb,mathtools}
\usepackage[scr=rsfs,cal=euler]{mathalfa}
\usepackage{xfrac}
\usepackage[unicode,hidelinks,bookmarks=false]{hyperref}
\newcommand{\ie}{i.e.\ }
\newcommand{\cf}{cf.\ }
\newcommand{\resp}{respectively\ }
\newcommand{\Subsection}{\subsection}
\providecommand{\nobreakdash}{\nobreak\hbox{-}}
\setlength{\emergencystretch}{2em}
\multlinegap0pt
\usepackage{multirow}
\usepackage{mathrsfs}
\usepackage{appendix}
\usepackage{textcomp}
\usepackage{manyfoot}
\usepackage{booktabs}
\usepackage{algorithm}
\usepackage{algorithmicx}
\usepackage{algpseudocode}
\usepackage{listings}
\newtheorem{theorem}{Theorem}
\newtheorem{proposition}[theorem]{Proposition}
\newtheorem{example}{Example}
\newtheorem{remark}{Remark}
\newtheorem{corollary}[theorem]{Corollary}
\newtheorem{definition}{Definition}
\newtheorem{conjecture}[theorem]{Conjecture}
\newtheorem{lemma}[theorem]{Lemma}
\begin{document}
\title[Permutation Polynomials over mathbbZ_n and T-quaternion Ring]{Permutation Polynomials over $\mathbb{Z}_{n}$ and T-quaternion Rings}
\author{A.Telveenus}
\date{}
\maketitle
\noindent\emph{Source and licence.} Permutation Polynomials over $\mathbb{Z}_{n}$ and T-quaternion Rings. Version arXiv:2608.09136v2 (CC BY 4.0). This material is distributed under the Creative Commons Attribution 4.0 International licence, http://creativecommons.org/licenses/by/4.0/, as recorded in the arXiv OAI licence metadata for this exact version and shown on the abstract page https://arxiv.org/abs/2608.09136v2. Full rights notice: licenses/arxiv-2608.09136-CC-BY-4.0.txt.\par\smallskip
\noindent\textit{Editorial scope.} Counted unit: the original \texttt{theorem} environment \texttt{rd25} at source lines 163--165 together with its complete proof at lines 166--188; the delivered document keeps source lines 160--194 so that every referenced label and every citation resolves inside it. Native editable source is the paper's own concatenated TeX; the AMS wrapper is editorial formatting only. No publisher class, upstream script or unrelated artwork is redistributed.
\section{Permutation Polynomials over $\mathbb{Z}n$}
For many values of $n$, the ring $\mathbb{Z}_{n}$ is not a field and it is a field only when $n$ is prime. In these cases, the zero divisors in $\mathbb{Z}_{n}$ play a crucial role in the construction of permutation polynomials.
A polynomial $f(X) \in \mathbb{Z}_{n}[X]$ is said to be a \textit{permutation polynomial} if $f:\mathbb{Z}_{n} \mapsto \mathbb{Z}_{n}$ is a bijection. If $f(X)+X$ is also a permutation polynomial, it is called a \textit{complete permutation polynomial}. 

\begin{theorem} \label{rd25}
If $n=p^{e}$ with prime $p>3$ and $e>1$, then there are $n \phi(n)$ linear, $[n-\phi(n)-1]n \phi(n)$ quadratic, and $[n-\phi(n)-1][n-\phi(n)]n \phi(n)$ cubic permutation polynomials in $\mathbb{Z}_{n}[X]$.
\end{theorem}

\begin{proof}
Case 1: Linear \\
Let $f(X)=AX+B$ be the linear permutation polynomial. The condition is $f(X_{1})=f(X_{2}) \quad \implies X_{1}=X_{2}$. Since the ring $\mathbb{Z}_{n}$ is finite, no need of checking the map $f$ is surjective. Thus we have $(X_{1}-X_{2})A=0$. So, $X_{1}=X_{2}$ is possible only when $A \in U(\mathbb{Z}_{n})$. This means that $A$ can take $\phi(n)$ and $B$ can take any values in $\mathbb{Z}_{n}$. Hence the result. \\
Case 2: Quadratic \\
Let $f(X)=AX^{2}+BX+C$ be the quadratic permutation polynomial. Proceding as in case 1, we get $X_{1}=X_{2}$ only when $A(X_{1}+X_{2})+B=2AX_{1}+B=u  \in U(\mathbb{Z}_{n})$ for any value of $X_{1}$ in 
$\mathbb{Z}_{n}$. Thus we have the following equations:
$$B=u$$
$$2A+B=u$$
$$4A+B=u$$
$$\cdots \cdots $$
$$2(p^{e}-1)A+B=u$$
From the above equations we get $B=u$ and $2A=4A=\cdots =2(p^{e}-1)A=0$. This means that $A$ is a zero divisor in $\mathbb{Z}_{n}$ with $A \neq 0$. Thus $A$ can take $[n-\phi(n)-1]$ values, $B$ can take $\phi(n)$ values being a unit, and $C$ can take any values in  $\mathbb{Z}_{n}$. Hence the result. \\
Case 3: Cubic \\
Let $f(X)=AX^{3}+BX^{2}+CX+D$ be the cubic permutation polynomial. Proceding as in case 1, we get $X_{1}=X_{2}$ only when $A(X_{1}^{2}+X_{1}X_{2}+X_{2}^{2})+B(X_{1}+X_{2})+C)=3AX_{1}^{2}+2BX_{1}+C=u  \in U(\mathbb{Z}_{n})$ for any value of $X_{1}$ in 
$\mathbb{Z}_{n}$. Thus we have the following equations:
$$C=u$$
$$3A+2B+C=u$$
$$12A+4B+C=u$$
$$\cdots \cdots $$
$$3(p^{e}-1)^{2}A+B(p^{e}-1)+C=u$$
The presence of $3$ and $2$ in the equations might affect the number of solutions if we take $n=2^{e}$ or $n=3^{e}$ as they are zero divisors in $n=2^{e}$ or $n=3^{e}$. This is the reason why we have excluded $p=2$ and $p=3$. \\
From the above equations we get $C=u$,  $3A=6A=\cdots =3(p^{e}-1)^{2}A=0$, and $2B=4B=6B \cdots = 2(p^{e}-1)B=0$. This means that $A$ is a zero divisor in $\mathbb{Z}_{n}$ with $A \neq 0$. Thus $A$ can take $[n-\phi(n)-1]$ values, $B$ is also a zero divisor and can take $[n-\phi(n)]$ values including zero,  $C$ being a unit it can take $\phi(n)$ values, and $D$ can take any values in  $\mathbb{Z}_{n}$. Hence the result. \\
\end{proof}

\begin{proposition}
If $p>3$ a prime, then the number of quadratic permutation polynomials is zero in $\mathbb{Z}_{p}[X]$. 
\end{proposition}
\begin{proof}
The propositon follows from the proof of case 2 of theorem \ref{rd25} as the only zero divisor is zero and A must take that value. This means that there are no quadratic permutation polynomials in $\mathbb{Z}_{p}[X]$ . 
\end{proof}

\end{document}
